% Practice Test 24 — Virginia Edition
% Questions: 30 (27 MC + 3 SA)
% Generated by scripts/generate_practice_tests.py

\begin{practiceQuestion}{1-4-q15}{mc}
\begin{questionText}
The point $(8, 52)$ is on a proportional graph. Which equation describes the relationship?
\end{questionText}
\choiceA{$y = 6x$}
\choiceB{$y = 6.5x$}
\choiceC{$y = 7x$}
\choiceD{$y = 8x$}
\correctAnswer{B}
\explanation{$k = \frac{52}{8} = 6.5$, so $y = 6.5x$.}
\end{practiceQuestion}

\begin{practiceQuestion}{2-1-q01}{mc}
\begin{questionText}
What is $25\%$ of $80$?
\end{questionText}
\choiceA{$15$}
\choiceB{$20$}
\choiceC{$25$}
\choiceD{$32$}
\correctAnswer{B}
\explanation{Convert to a decimal: $25\% = 0.25$. Multiply: $0.25 \times 80 = 20$.}
\end{practiceQuestion}

\begin{practiceQuestion}{2-2-q05}{mc}
\begin{questionText}
Solve using a proportion: $72$ is $90\%$ of what number?
\end{questionText}
\choiceA{$64.8$}
\choiceB{$75$}
\choiceC{$78$}
\choiceD{$80$}
\correctAnswer{D}
\explanation{$\frac{72}{w} = \frac{90}{100}$. Cross-multiply: $90w = 7200$. Divide: $w = 80$.}
\end{practiceQuestion}

\begin{practiceQuestion}{2-3-q06}{mc}
\begin{questionText}
What multiplier represents a $40\%$ decrease?
\end{questionText}
\choiceA{$0.40$}
\choiceB{$0.60$}
\choiceC{$1.40$}
\choiceD{$1.60$}
\correctAnswer{B}
\explanation{A $40\%$ decrease means you keep $100\% - 40\% = 60\%$. The multiplier is $0.60$.}
\end{practiceQuestion}

\begin{practiceQuestion}{2-7-q09}{mc}
\begin{questionText}
If your estimate is exactly correct, what is the percent error?
\end{questionText}
\choiceA{$0\%$}
\choiceB{$1\%$}
\choiceC{$50\%$}
\choiceD{$100\%$}
\correctAnswer{A}
\explanation{If estimated $=$ actual, then $|estimated - actual| = 0$, so percent error $= 0\%$.}
\end{practiceQuestion}

\begin{practiceQuestion}{3-2-q07}{mc}
\begin{questionText}
What is $(-3) + (-9) + 4$?
\end{questionText}
\choiceA{$-8$}
\choiceB{$-16$}
\choiceC{$10$}
\choiceD{$2$}
\correctAnswer{A}
\explanation{Add the negatives first: $(-3) + (-9) = -12$. Then $(-12) + 4 = -8$.}
\end{practiceQuestion}

\begin{practiceQuestion}{3-3-q11}{mc}
\begin{questionText}
A scuba diver is at $-25$ feet. She swims up $10$ feet then descends $7$ feet. What is her new depth?
\end{questionText}
\choiceA{$-22$ feet}
\choiceB{$-32$ feet}
\choiceC{$-42$ feet}
\choiceD{$-12$ feet}
\correctAnswer{A}
\explanation{$-25 + 10 = -15$. Then $-15 - 7 = -15 + (-7) = -22$ feet.}
\end{practiceQuestion}

\begin{practiceQuestion}{3-4-q30}{gsa}
\begin{questionText}
A thermometer shows these temperature readings over three recordings.

\begin{center}
\begin{tikzpicture}
  \draw[thick,->] (0,-4.5) -- (0,4.5);
  \foreach \y in {-4,...,4} {
    \draw (-0.15,\y) -- (0.15,\y);
    \node[right=4pt] at (0,\y) {$\y°$C};
  }
  \fill[funBlue] (0,-2.5) circle (4pt);
  \node[left=10pt] at (0,-2.5) {6 AM: $-2.5°$C};
  \draw[thick,->,funGreen] (-0.7,-2.5) -- (-0.7,1.3);
  \node[left] at (-0.7,-0.6) {\footnotesize $+3.8°$};
  \draw[thick,->,funRed] (-1.4,1.3) -- (-1.4,-0.4);
  \node[left] at (-1.4,0.45) {\footnotesize $-1.7°$};
\end{tikzpicture}
\end{center}

What is the final temperature reading?
\end{questionText}
\correctAnswer{$-0.4°$C}
\explanation{Start at $-2.5$. Rise $3.8$: $-2.5 + 3.8 = 1.3°$C. Drop $1.7$: $1.3 - 1.7 = -0.4°$C.}
\end{practiceQuestion}

\begin{practiceQuestion}{4-3-q07}{mc}
\begin{questionText}
A student expands $-3(x - 4)$ and writes $-3x - 12$. What mistake did the student make?
\end{questionText}
\choiceA{Forgot to distribute to the second term}
\choiceB{Did not change the sign when multiplying $-3 \times (-4)$}
\choiceC{Multiplied only the first term}
\choiceD{Added instead of multiplying}
\correctAnswer{B}
\explanation{$-3 \times (-4)$ should equal $+12$, not $-12$. The correct expansion is $-3x + 12$.}
\end{practiceQuestion}

\begin{practiceQuestion}{4-4-q29}{gsa}
\begin{questionText}
The total area of the two rectangles shown below is $8x + 12$. Both rectangles have the same height $h$. Find the value of $h$.

\begin{center}
\begin{tikzpicture}[scale=0.7]
  \draw[thick, fill=funBlue!15] (0,0) rectangle (4,2.5);
  \node[below] at (2,0) {$2x$};
  \node[left] at (0,1.25) {$h$};
  \node at (2,1.25) {};
  \draw[thick, fill=funGreen!15] (5,0) rectangle (8,2.5);
  \node[below] at (6.5,0) {$3$};
  \node[right] at (8,1.25) {$h$};
\end{tikzpicture}
\end{center}
\end{questionText}
\correctAnswer{$h = 4$}
\explanation{Total area $= h(2x) + h(3) = h(2x + 3) = 8x + 12$. Factor: $8x + 12 = 4(2x + 3)$. So $h = 4$.}
\end{practiceQuestion}

\begin{practiceQuestion}{5-4-q10}{mc}
\begin{questionText}
Solve $\dfrac{x}{6} - \dfrac{2}{3} = \dfrac{1}{2}$.
\end{questionText}
\choiceA{$x = 3$}
\choiceB{$x = 7$}
\choiceC{$x = -1$}
\choiceD{$x = 10$}
\correctAnswer{B}
\explanation{Multiply every term by $6$: $x - 4 = 3$. Add $4$: $x = 7$. Check: $\frac{7}{6} - \frac{2}{3} = \frac{7}{6} - \frac{4}{6} = \frac{3}{6} = \frac{1}{2}$ \checkmark.}
\end{practiceQuestion}

\begin{practiceQuestion}{5-5-q01}{mc}
\begin{questionText}
Solve $2x + 3 > 11$.
\end{questionText}
\choiceA{$x > 4$}
\choiceB{$x > 7$}
\choiceC{$x < 4$}
\choiceD{$x > 14$}
\correctAnswer{A}
\explanation{Subtract $3$: $2x > 8$. Divide by $2$: $x > 4$.}
\end{practiceQuestion}

\begin{practiceQuestion}{5-6-q13}{mc}
\begin{questionText}
True or false: The graph of $x < 5$ and $x \leq 5$ look exactly the same.
\end{questionText}
\choiceA{True — both shade to the left of $5$}
\choiceB{False — $x < 5$ has an open circle, $x \leq 5$ has a closed circle}
\choiceC{False — they shade in different directions}
\choiceD{True — the circle type does not matter}
\correctAnswer{B}
\explanation{$x < 5$ uses an open circle (not including $5$) while $x \leq 5$ uses a closed circle (including $5$). Both shade left.}
\end{practiceQuestion}

\begin{practiceQuestion}{6-1-q02}{mc}
\begin{questionText}
A blueprint uses the scale $1$ in $= 6$ ft. A hallway is $3.5$ in long on the blueprint. What is the actual length?
\end{questionText}
\choiceA{$18$ ft}
\choiceB{$21$ ft}
\choiceC{$24$ ft}
\choiceD{$9.5$ ft}
\correctAnswer{B}
\explanation{$3.5 \times 6 = 21$ ft.}
\end{practiceQuestion}

\begin{practiceQuestion}{6-2-q16}{mc}
\begin{questionText}
A model uses scale $1:25$. A new model uses scale $1:75$. A window that is $6$ cm on the first model becomes how long on the second?
\end{questionText}
\choiceA{$2$ cm}
\choiceB{$6$ cm}
\choiceC{$18$ cm}
\choiceD{$150$ cm}
\correctAnswer{A}
\explanation{Scale ratio: $\frac{25}{75} = \frac{1}{3}$. New length: $6 \times \frac{1}{3} = 2$ cm.}
\end{practiceQuestion}

\begin{practiceQuestion}{6-5-q10}{mc}
\begin{questionText}
A cube can be sliced to produce a hexagonal cross-section. How?
\end{questionText}
\choiceA{By cutting parallel to the base}
\choiceB{By cutting diagonally through all six faces}
\choiceC{By cutting vertically through the center}
\choiceD{It is impossible to get a hexagon from a cube}
\correctAnswer{B}
\explanation{A diagonal plane that passes through all six faces of a cube creates a regular hexagonal cross-section.}
\end{practiceQuestion}

\begin{practiceQuestion}{6-6-q07}{mc}
\begin{questionText}
Two vertical angles are $(5x - 20)°$ and $(3x + 10)°$. What is the measure of each angle?
\end{questionText}
\choiceA{$45°$}
\choiceB{$55°$}
\choiceC{$60°$}
\choiceD{$75°$}
\correctAnswer{B}
\explanation{Vertical angles are equal: $5x - 20 = 3x + 10$. Subtract $3x$: $2x - 20 = 10$. Add $20$: $2x = 30$. Divide: $x = 15$. Angle: $5(15) - 20 = 55°$.}
\end{practiceQuestion}

\begin{practiceQuestion}{6-7-q04}{mc}
\begin{questionText}
Point $D(2, 6)$ is rotated $90°$ clockwise about the origin. What are the coordinates of $D'$?
\end{questionText}
\choiceA{$(-6, 2)$}
\choiceB{$(6, -2)$}
\choiceC{$(-2, -6)$}
\choiceD{$(6, 2)$}
\correctAnswer{B}
\explanation{$90°$ clockwise: $(x, y) \to (y, -x)$. $(2, 6) \to (6, -2)$.}
\end{practiceQuestion}

\begin{practiceQuestion}{6-8-q09}{mc}
\begin{questionText}
Two similar triangles have corresponding sides of $7$ and $21$. What is the ratio of their areas?
\end{questionText}
\choiceA{$1:3$}
\choiceB{$1:7$}
\choiceC{$1:9$}
\choiceD{$7:21$}
\correctAnswer{C}
\explanation{Scale factor: $\frac{21}{7} = 3$. Area ratio: $3^2 = 9$. So the ratio is $1:9$.}
\end{practiceQuestion}

\begin{practiceQuestion}{7-4-q15}{mc}
\begin{questionText}
A shape is made of a rectangle $8$ in by $6$ in and two semicircles, each with a diameter of $6$ in, attached to both short sides. What is the total area? Use $\pi \approx 3.14$.
\end{questionText}
\choiceA{$48$ in$^2$}
\choiceB{$76.26$ in$^2$}
\choiceC{$104.52$ in$^2$}
\choiceD{$133.04$ in$^2$}
\correctAnswer{B}
\explanation{Rectangle: $8 \times 6 = 48$ in$^2$. Two semicircles $=$ one full circle with $r = 3$: $3.14 \times 9 = 28.26$ in$^2$. Total: $48 + 28.26 = 76.26$ in$^2$.}
\end{practiceQuestion}

\begin{practiceQuestion}{7-5-q18}{mc}
\begin{questionText}
A rectangular prism has a surface area of $94$ cm$^2$. Its length is $5$ cm and width is $3$ cm. What is its height?
\end{questionText}
\choiceA{$2$ cm}
\choiceB{$3$ cm}
\choiceC{$4$ cm}
\choiceD{$5$ cm}
\correctAnswer{C}
\explanation{$94 = 2(5 \times 3) + 2(5 \times h) + 2(3 \times h) = 30 + 10h + 6h = 30 + 16h$. So $16h = 64$, $h = 4$ cm.}
\end{practiceQuestion}

\begin{practiceQuestion}{7-6-q01}{mc}
\begin{questionText}
What is the volume of a rectangular prism with length $6$ cm, width $4$ cm, and height $3$ cm?
\end{questionText}
\choiceA{$13$ cm$^3$}
\choiceB{$24$ cm$^3$}
\choiceC{$48$ cm$^3$}
\choiceD{$72$ cm$^3$}
\correctAnswer{D}
\explanation{$V = lwh = 6 \times 4 \times 3 = 72$ cm$^3$.}
\end{practiceQuestion}

\begin{practiceQuestion}{7-7-q10}{mc}
\begin{questionText}
A cylinder has a diameter of $10$ in and a height of $6$ in. What is its volume? Use $\pi \approx 3.14$.
\end{questionText}
\choiceA{$188.4$ in$^3$}
\choiceB{$471$ in$^3$}
\choiceC{$942$ in$^3$}
\choiceD{$1{,}884$ in$^3$}
\correctAnswer{B}
\explanation{$r = 5$ in. $V = \pi r^2 h = 3.14 \times 25 \times 6 = 471$ in$^3$.}
\end{practiceQuestion}

\begin{practiceQuestion}{8-1-q02}{mc}
\begin{questionText}
A researcher wants to know the average height of seventh graders in a state. She measures the heights of $100$ seventh graders from different schools. What is the sample?
\end{questionText}
\choiceA{All seventh graders in the state}
\choiceB{All students in the schools she visited}
\choiceC{The $100$ seventh graders she measured}
\choiceD{The tallest seventh graders}
\correctAnswer{C}
\explanation{The sample is the smaller group chosen from the population — the $100$ students measured.}
\end{practiceQuestion}

\begin{practiceQuestion}{8-3-q26}{gmc}
\begin{questionText}
The box plots below compare test scores for Class A and Class B.

\begin{center}
\begin{tikzpicture}[scale=0.1]
  \draw[thick,->] (49,0) -- (101,0);
  \foreach \x in {50,55,...,100} {
    \draw (\x,0.5) -- (\x,-0.5);
    \node[below,font=\small] at (\x,-0.5) {$\x$};
  }
  % Class A: min=60, Q1=70, med=78, Q3=85, max=95
  \draw[thick] (60,6) -- (95,6);
  \draw[thick,fill=funBlue!30] (70,4) rectangle (85,8);
  \draw[very thick,funBlue] (78,4) -- (78,8);
  \draw[thick] (60,4) -- (60,8);
  \draw[thick] (95,4) -- (95,8);
  \node[right,font=\small] at (96,6) {Class A};
  % Class B: min=50, Q1=58, med=65, Q3=72, max=80
  \draw[thick] (50,14) -- (80,14);
  \draw[thick,fill=funGreen!30] (58,12) rectangle (72,16);
  \draw[very thick,funGreen] (65,12) -- (65,16);
  \draw[thick] (50,12) -- (50,16);
  \draw[thick] (80,12) -- (80,16);
  \node[right,font=\small] at (81,14) {Class B};
\end{tikzpicture}
\end{center}

Which statement is best supported by the box plots?
\end{questionText}
\choiceA{Class B scored higher than Class A}
\choiceB{Class A has a higher median and a larger IQR}
\choiceC{Class A has a higher median and both classes have the same IQR}
\choiceD{Class A and Class B overlap, but Class A's median is higher}
\correctAnswer{D}
\explanation{Class A median $= 78$, Class B median $= 65$. Class A's range ($60$--$95$) overlaps with Class B's range ($50$--$80$), but Class A is higher overall.}
\end{practiceQuestion}

\begin{practiceQuestion}{8-4-q16}{mc}
\begin{questionText}
Group A: $\{60, 65, 70, 75, 80\}$. Group B: $\{40, 55, 70, 85, 100\}$. Both have a mean of $70$. Which group is more variable?
\end{questionText}
\choiceA{Group A}
\choiceB{Group B}
\choiceC{They are equally variable}
\choiceD{Cannot be determined}
\correctAnswer{B}
\explanation{Group A range $= 20$, Group B range $= 60$. Group B's data is far more spread out.}
\end{practiceQuestion}

\begin{practiceQuestion}{8-5-q27}{gmc}
\begin{questionText}
The back-to-back stem-and-leaf plot compares quiz scores for two classes.

\begin{center}
\begin{tabular}{r|c|l}
\textbf{Class A} & \textbf{Stem} & \textbf{Class B} \\
\hline
$5\;\;2$ & $6$ & $3\;\;5\;\;8$ \\
$8\;\;5\;\;3\;\;0$ & $7$ & $2\;\;4\;\;9$ \\
$6\;\;4$ & $8$ & $1\;\;3\;\;5\;\;7$ \\
$2$ & $9$ & $0\;\;5$
\end{tabular}

Key: $5 \mid 7$ means $75$ (Class A) and $7 \mid 2$ means $72$ (Class B)
\end{center}

Which class has a higher median?
\end{questionText}
\choiceA{Class A}
\choiceB{Class B}
\choiceC{They have the same median}
\choiceD{Cannot be determined}
\correctAnswer{B}
\explanation{Class A data ($9$ values): $62, 65, 70, 73, 75, 78, 84, 86, 92$. Median $= 75$. Class B data ($12$ values): $63, 65, 68, 72, 74, 79, 81, 83, 85, 87, 90, 95$. Median $= \frac{79+81}{2} = 80$. Class B has a higher median.}
\end{practiceQuestion}

\begin{practiceQuestion}{9-3-q09}{mc}
\begin{questionText}
A spinner with $4$ equal sections is spun $80$ times. Section A appears $25$ times. How does the experimental probability compare to the theoretical probability?
\end{questionText}
\choiceA{Experimental ($0.3125$) is higher than theoretical ($0.25$)}
\choiceB{Experimental ($0.25$) equals theoretical ($0.25$)}
\choiceC{Experimental ($0.20$) is lower than theoretical ($0.25$)}
\choiceD{Experimental ($0.3125$) is lower than theoretical ($0.50$)}
\correctAnswer{A}
\explanation{Experimental: $\frac{25}{80} = 0.3125$. Theoretical: $\frac{1}{4} = 0.25$. The experimental is slightly higher.}
\end{practiceQuestion}

\begin{practiceQuestion}{9-6-q18}{mc}
\begin{questionText}
Three coins are flipped. What is the probability of getting NO heads (all tails)?
\end{questionText}
\choiceA{$\frac{1}{2}$}
\choiceB{$\frac{1}{4}$}
\choiceC{$\frac{1}{8}$}
\choiceD{$\frac{3}{8}$}
\correctAnswer{C}
\explanation{Only one outcome out of $8$ is all tails (TTT). $P = \frac{1}{8}$.}
\end{practiceQuestion}

\begin{practiceQuestion}{9-7-q24}{sa}
\begin{questionText}
Describe how to use a standard die to simulate a $\frac{1}{3}$ probability event.
\end{questionText}
\correctAnswer{Roll the die; let $1$ and $2$ represent the event (or any $2$ numbers out of $6$)}
\explanation{Choose $2$ of the $6$ faces to represent the event. $P = \frac{2}{6} = \frac{1}{3}$.}
\end{practiceQuestion}
